% 第14章 评测污染审计
\section{评测数据污染与不可比设定审计}\label{sec:audit}

这一章是整份综述里最不讨喜、但可能最有用的一部分。金融机器学习的论文浩如烟海，真正值得重读的却不多，原因不在于方法不够新，而在于\textbf{评测本身出了问题}。审计围绕五类风险展开，前四类是老问题在 AI 时代的放大，第五类是 AI 带来的新问题。

\subsection{幸存者偏差与上市偏差}

用“今天还在交易的股票”回溯过去，会系统性高估收益，因为退市、被并购、破产的样本被剔除了。含完整退市记录的数据库（如 CRSP）能把这一类偏差压到较低水平；而用现成分股名单回测的做法，在大量“深度选股”论文里仍然常见。判断一条结论是否可信，第一件事就是看它用的是时点（point-in-time）成分股，还是事后名单。

\subsection{前瞻偏差（look-ahead）}

前瞻偏差在 AI 交易里无处不在，且形态隐蔽：

\begin{itemize}
  \item \textbf{特征与标签的时点错位}：$t$ 期收盘算出信号，却假设 $t$ 期开盘成交。
  \item \textbf{财报的可用日}：用报告期而非实际披露日对齐财务数据。
  \item \textbf{标准化泄漏}：用全样本均值方差做特征标准化，把未来信息带进训练。
  \item \textbf{数据处理的全样本操作}：去极值、中性化、缺失值填充若在全样本上做，同样泄漏。
  \item \textbf{图结构的时点}：第 \ref{sec:gnn} 章提到，用“今天的行业/供应链关系”回测十年前。
\end{itemize}

这些错误不会让代码报错，只会让回测好看，因此极难自查。工程上唯一可靠的防线是\textbf{用“事件时间线”重放数据，并在每个时点只允许访问当时可见的信息}。

\subsection{数据窥探与多重检验}

这是量化研究最古老的陷阱：测一万条规则，总能找到一条赚钱的。White（2000）\cite{white2000} 的 Reality Check 用 bootstrap 对大量候补规则做联合检验；Sullivan、Timmermann 与 White（1999）\cite{sullivan1999} 对约近八千条技术规则做了联合检验，结论是整个规则集合并不显著优于“什么都不做”。Harvey、Liu 与 Zhu（2016）\cite{harveylz2016} 给出截面研究的修正：在数百个已发表因子面前，新因子的 $t$ 阈值应从 2 提到约 3。

AI 把这个问题放大了一个量级。神经网络的超参数（层数、宽度、学习率、正则、早停点）、特征工程、甚至随机种子，都构成隐含的“试验次数”。如果模型选择是在同一个样本外集合上反复进行，这个“样本外”就已经被污染，不再是真正的样本外。一个常被忽略的事实是：\textbf{深度学习选股论文报告的“样本外”夏普，多数是在被反复用来选模型的那个集合上算出来的。}需要补充的是，复现危机并非只存在于股票因子——近期对债券因子的复现研究\cite{bondreplication2026} 显示，同类问题在其他资产类别上同样严重。

\subsection{回测过拟合与选择性报告}

Bailey 等（2014）\cite{bailey2014pseudo} 用“伪数学”一词批评了回测报告中的常见误导；Bailey 与 L\'opez de Prado（2014）\cite{bailey2014dsr} 提出\textbf{缩水夏普比率}（DSR），把试验次数、样本长度与收益非正态都纳入显著性；Bailey 等（2016）\cite{bailey2016pbo} 提出\textbf{回测过拟合概率}（PBO）与组合对称交叉验证（CSCV）。L\'opez de Prado（2018）\cite{lopezdeprado2018} 系统整理了金融机器学习的评测纪律。

这里有一条对 AI 交易尤其重要的量化关系：\textbf{若在 $M$ 个候选策略中选样本内最优者，其样本内夏普的期望最大值随 $\sqrt{\ln M}$ 增长，即使所有策略都是纯噪声。}这意味着“多试几次”本身就会制造出看起来显著的策略——而强化学习与大规模超参数搜索，恰恰是“多试几次”的极端形式。

\subsection{AI 时代的新污染：模型知道答案}

前三类是老问题，第四类是第五类的基础，而这一节要讲的是 AI 带来的、可能最严重的\textbf{知识泄漏}。

\textbf{（一）大模型的预训练语料泄漏。}用一个 2024 年发布的 LLM 回测 2020--2023 年的新闻，如果这段文本在其预训练语料中，模型是在“背答案”而非“推理”。第 \ref{sec:nlp} 章已讨论过这一点，这里再强调其后果：\textbf{这类泄漏不会表现为异常高的精度，而会表现为“看起来合理的高收益”，因此极难通过常规手段发现。}唯一可靠的处理是使用模型知识截止日之后的前向样本。

\textbf{（二）生成数据中的训练/测试泄漏。}用生成模型造数据来增强训练集时，若生成器的训练数据与测试集重叠，测试集就被间接泄漏。合成数据的“信息回收”是隐蔽的第二条泄漏通道。

\textbf{（三）强化学习的环境泄漏。}用历史回放作为 RL 环境时，测试期的行情出现在训练环境里，策略对“未来”提前适应。

\subsection{跨论文不可比设定}

即便单篇论文内部干净，把多篇论文放在一起比较仍然困难，因为设定差异巨大：

\begin{table}[htbp]
\centering
\caption{跨论文不可比设定的典型来源}
\label{tab:incomparable}
\begin{tabular}{@{}>{\raggedright\arraybackslash}p{3.2cm}>{\raggedright\arraybackslash}p{11.4cm}@{}}
\toprule
设定维度 & 差异举例 \\
\midrule
任务定义 & “预测次日方向”“预测 5 日收益”“排序”是三种不可比的任务，指标也不同 \\
样本区间 & 含 2020 年疫情行情的样本与不含的样本，不可比 \\
资产池 & 美股大盘、A 股全市场、加密永续合约的微结构差异巨大 \\
频率 & 日频选股、分钟级执行、逐笔做市，对模型的要求完全不同 \\
成本口径 & 零成本、固定费率、含冲击成本，三者可差一个数量级 \\
标签构造 & 未来 1 期 vs $k$ 期、绝对收益 vs 超额收益，信号方向可能相反 \\
基准选择 & 与“无脑买入持有”比，还是与“最强已知因子”比，结论不同 \\
评测协议 & 随机切分 vs 滚动前向；单次切分 vs walk-forward \\
\bottomrule
\end{tabular}
\end{table}

\subsection{一个最低可接受的评测协议}

把上述五类风险合起来，可以给出一份“最低可接受”的评测协议。它不是理想标准，而是\textbf{低于这个标准的结果就不该被当作证据}的下限：

\begin{enumerate}
  \item \textbf{数据}：声明数据来源、样本区间、时点对齐方式；公布或描述成分股/关系/文本的时间戳口径。
  \item \textbf{切分}：使用滚动前向或明确的前向留出段；禁止用测试段调参；随机切分一律视为不合格。
  \item \textbf{成本}：至少报告固定费率与含冲击两档；说明冲击模型与参与率上限。
  \item \textbf{基准}：与“最强简单基线”对照（买入持有、经典因子、VWAP 等），并说明基线的口径。
  \item \textbf{显著性}：报告试验次数 $N$、样本长度 $T$，并给出 DSR 与 PBO；“未报告试验次数”按不合格处理。
  \item \textbf{泄漏}：若使用 LLM 或生成式模型，必须证明测试样本不在其训练语料/训练分布的知识范围内。
  \item \textbf{可复现}：给出代码或可重建的完整流程说明；数据私有者需说明替代验证途径（如前向样本、外部数据集复现）。
\end{enumerate}

这份协议里，前四条是“统计上诚实”，后三条是“AI 时代新增的要求”。本文认为，未来一段时间内，领域最大的边际改进空间不在算法，而在于把这份协议变成投稿与发布的默认门槛——就像临床试验的 CONSORT 之于医学那样。

\subsection{审计结论}

把上面的风险合起来，本文给出一条审计判据，供读者在读到任何一条“AI 交易”结论时套用：

\begin{enumerate}
  \item 是否用时点数据（成分股、财报、关系、文本）？
  \item 是否报告了扣成本后的样本外结果？
  \item 是否报告了试验次数并做多重检验校正（DSR/PBO）？
  \item 若用到 LLM，是否证明了模型知识截止日早于样本期？
  \item 是否与强基线（而非“无脑持有”）对照？
  \item 是否提供了可复现的代码与数据口径？
\end{enumerate}

\textbf{六问全过，结论才值得认真对待；六问中间任何一问为否，都应在结论上打一个显著的折扣。}本文的证据表与复现状态核验表，正是按这套判据逐条标注的。本章要记住的一句话是：\textbf{金融机器学习里，最贵的错误不是模型选错了，而是你用来选模型的那把尺子本身就是歪的。}
